\subsubsection{Proof of Theorem~\ref{thm:dirichlet-continuous-data-existence}}\label{subsec:perron-existence-proof}

We are now ready to prove our main result. Let $S_g$ be the set of all subsolutions to the Dirichlet problem, and let
\begin{equation}\label{eq:perron-supremum-of-subsolutions}
u(\bm{x})=\sup_{v\in S_g}v(\bm{x}).
\end{equation}
Since every subsolution is less than or equal to every supersolution and both subsolutions and supersolutions exist, the supremum in \eqref{eq:perron-supremum-of-subsolutions} is finite, and $u$ is well defined.

\begin{Lemma}
\label{lem:perron-envelope-harmonic}
The function $u$ is harmonic in $\Omega$.
\end{Lemma}

\begin{Proof}
Let $\bm{x}$ be an arbitrary point in $\Omega$ and let $v_m$ be a sequence in $S_g$ such that $v_m(\bm{x})\to u(\bm{x})$. Clearly $v_m$ is bounded from above (by any supersolution), we may also assume it is bounded from below; if necessary, replace $v_m$ by $\max(v_m,v_0)$, where $v_0$ is any subsolution. Choose $R$ such that the closure of $B=B_R(\bm{x})$ is contained in $\Omega$ and let $V_m$ be the harmonic lifting of $v_m$ with respect to $B$. Then $V_m(\bm{x})\to u(\bm{x})$ and, by Theorem~\ref{thm:harmonic-sequence-compactness}, $V_m$ has a subsequence $V_{m_k}$ which converges on $B$ to a harmonic function $v$; the convergence is uniform on every compact subset of $B$. Clearly $v\leq u$ in $B$. We shall now show that in fact $v=u$ in $B$. This clearly implies the lemma.

Assume on the contrary that $v(\bm{y})<u(\bm{y})$ for some $\bm{y}\in B$. Then there exists a function $W\in S_g$ such that $v(\bm{y})<W(\bm{y})$. Let $w_k=\max(W,V_{m_k})$ and let $W_k$ be the harmonic lifting of $w_k$ with respect to $B$. As before, a subsequence of $W_k$ converges to a function $w$, harmonic in $B$. Then clearly $v\leq w$ on $B$ and $v(\bm{x})=w(\bm{x})$. By the maximum principle, this implies $v=w$ on $B$, contradicting the choice of $W$.
\end{Proof}

It remains to be shown that $u$ is actually continuous up to $\partial\Omega$ and assumes the given boundary data. Let $\bm{\xi}$ be a point on $\partial\Omega$. Since $\partial\Omega$ was assumed of class $C^2$, there is a ball $B=B_R(\bm{y})$ in the exterior of $\Omega$ such that $\overline B\cap\overline\Omega=\{\bm{\xi}\}$. We now define
\begin{equation}\label{eq:perron-exterior-ball-barrier}
w(\bm{x})=\begin{cases}R^{2-n}-|\bm{x}-\bm{y}|^{2-n},&n\geq3\\\log\dfrac{|\bm{x}-\bm{y}|}{R},&n=2.\end{cases}
\end{equation}
Then
\begin{enumerate}
\item $w$ is superharmonic (actually harmonic) in $\Omega$.
\item $w>0$ on $\overline\Omega\setminus\{\bm{\xi}\}$, $w(\bm{\xi})=0$.
\end{enumerate}

A continuous function with properties (1) and (2) is called a \emph{barrier} at $\bm{\xi}$ relative to $\Omega$. We have the following result.

\begin{Lemma}
\label{lem:perron-boundary-continuity}
Let $u$ be the harmonic function in $\Omega$ constructed above and let $\bm{\xi}\in\partial\Omega$. Then $u(\bm{x})\to g(\bm{\xi})$ as $\bm{x}\to\bm{\xi}$.
\end{Lemma}

\begin{Proof}
Choose $\epsilon>0$ and let $M=\max_{\partial\Omega}|g|$. Let $w$ be a barrier at $\bm{\xi}$ and let $\delta$ and $k$ be such that $|g(\bm{x})-g(\bm{\xi})|<\epsilon$ for $\bm{x}\in\partial\Omega$, $|\bm{x}-\bm{\xi}|<\delta$ and $kw(\bm{x})\geq2M$ for $\bm{x}\in\Omega$, $|\bm{x}-\bm{\xi}|\geq\delta$. The functions $g(\bm{\xi})+\epsilon+kw(\bm{x})$ and $g(\bm{\xi})-\epsilon-kw(\bm{x})$ are, respectively, a supersolution and subsolution, and hence
\begin{equation}\label{eq:perron-barrier-boundary-comparison}
g(\bm{\xi})-\epsilon-kw(\bm{x})\leq u(\bm{x})\leq g(\bm{\xi})+\epsilon+kw(\bm{x}).
\end{equation}
Since $w(\bm{x})\to0$ as $\bm{x}\to\bm{\xi}$, the lemma is immediate.
\end{Proof}

It is clear from the proof that the smoothness assumption on $\partial\Omega$ can be relaxed as long as barriers can be constructed. The above construction works for all convex domains. In other cases, alternative barriers are often available, see, e.g., Problem~\ref{ex:planar-local-logarithmic-barrier}.

\paragraph*{Problems}

\begin{enumerate}[label=\thesection.\arabic*.,ref=\thesection.\arabic*,start=8,series=book-problems]

\item\label{ex:poisson-formula-harmonicity} Verify that the function given by \eqref{eq:poisson-ball-solution-formula} is harmonic.

\item\label{ex:subharmonic-harmonic-lifting} Prove the second claim in the subsection on subharmonic functions.

\item\label{ex:smooth-subharmonic-laplacian} Show that if $u$ is of class $C^2$ and subharmonic in the sense defined in this section, then $\Delta u\geq0$.

\item\label{ex:compact-equicontinuity-uniform} Let the sequence $f_m$ be equicontinuous at each point of a compact set $S$. Show that it is uniformly equicontinuous on $S$.

\item\label{ex:punctured-disk-dirichlet-nonexistence} Let $\Omega=\{(x,y)\in\mathbb{R}^2\mid0<x^2+y^2<1\}$. Prove that there is no solution to the Dirichlet problem $\Delta u=0$ in $\Omega$, $u(\bm{x})=1$ for $x^2+y^2=1$, $u(0)=0$. Hint: Show first that if there is a solution, then there is also a radially symmetric solution.

\item\label{ex:local-barrier-to-global-barrier} A function $w$ is called a local barrier at $\bm{\xi}$ relative to $\Omega$ if the conditions (i) and (ii) in the definition hold only in some neighborhood $N$ of $\bm{\xi}$. Show that if there is a local barrier, there is a barrier. Hint: Let $B$ be a small ball centered at $\bm{\xi}$ and $m=\inf_{N\setminus B}w$. Define $\bar w=\min(m,w)$ on $B$ and $\bar w=m$ otherwise.

\item\label{ex:planar-local-logarithmic-barrier} Let $\Omega$ be a domain in $\mathbb{R}^2$ with the origin on its boundary. Let $r$ and $\theta$ denote polar coordinates. Assume that a single-valued branch of $\theta$ exists in $\Omega\cap N$, where $N$ is a neighborhood of the origin. Show that a local barrier exists at the origin. Hint: Consider the analytic function $1/\log z$.

\item\label{ex:perron-proof-unbounded-domain} Why does the proof of Theorem~\ref{thm:dirichlet-continuous-data-existence} not work for unbounded domains?

\item\label{ex:unit-ball-volume-gamma-function} Show that $\omega_n=\pi^{n/2}/\Gamma(\frac n2+1)$. Hint: Evaluate \allowbreak $\int_{\mathbb{R}^n}\exp(-|\bm{x}|^2)\,\mathrm{d}\bm{x}$ in both Cartesian and polar coordinates. By comparing the two expressions, you find a formula for $\Omega_n=n\omega_n$, the surface area of the unit sphere.

\end{enumerate}
